\documentclass[11pt]{article}
\usepackage[utf8]{inputenc}
\usepackage[T1]{fontenc}
\usepackage{lmodern}
\usepackage{amsmath,amssymb}
\usepackage{graphicx}
\usepackage{booktabs}
\usepackage{ulem}
\usepackage{fixltx2e}
\usepackage[margin=20mm,a4paper]{geometry}
\usepackage{hyperref}
\hypersetup{colorlinks=true,linkcolor=blue,urlcolor=blue}
\usepackage{listings}
\lstset{basicstyle=\ttfamily\small,breaklines=true}
\usepackage{kotex}

\date{}

\begin{document}


\section*{My Document Converter 샘플 문서}\label{my-document-converter-샘플-문서}

\section*{소개 (Introduction)}\label{소개-introduction}

이 문서는 \textbf{My Document Converter} 가 지원하는 모든 출력 형식의 샘플을 만드는 데 쓰이는
원본입니다. \emph{강조}, \textbf{굵게}, \sout{취소선}, \verb|인라인 코드|,
H\textsubscript{2}O 와 x\textsuperscript{2}, 그리고
\href{https://example.com}{링크}를 담고 있습니다. This paragraph is in English so
that Latin-only formats still show something readable.

\subsection*{목록 (Lists)}\label{목록-lists}

\begin{itemize}
\item 첫 번째 항목
\item 두 번째 항목
\begin{itemize}
\item 중첩된 항목
\item 또 하나의 중첩 항목
\end{itemize}
\item 세 번째 항목
\end{itemize}

\begin{enumerate}
\item 순서가 있는 항목
\item 두 번째
\item 세 번째
\end{enumerate}

\begin{itemize}
\item ☒ 끝난 일
\item ☐ 남은 일
\end{itemize}

; 용어
: 용어에 대한 정의입니다.
; Pandoc
: 범용 문서 변환기. 이 프로그램은 같은 형식 집합을 자체 엔진으로 다룹니다.

\subsection*{코드 (Code)}\label{코드-code}

\begin{lstlisting}[language=js]
function greet(name) {
  console.log(`Hello, ${name}!`)
}
greet('world')
\end{lstlisting}

\subsection*{표 (Table)}\label{표-table}

\begin{table}[htbp]
\centering
\begin{tabular}{llll}
\toprule
형식 & 확장자 & 읽기 & 쓰기 \\
\midrule
Markdown & .md & O & O \\
HTML & .html & O & O \\
DOCX & .docx & O & O \\
EPUB & .epub & O & O \\
\bottomrule
\end{tabular}
\end{table}

\subsection*{인용과 수식 (Quote and math)}\label{인용과-수식-quote-and-math}

\begin{quote}
좋은 문서는 어떤 형식으로도 읽힙니다. --- 누군가
\end{quote}

인라인 수식 <latex>E = mc\^{}2</latex> 과 블록 수식:

<latex>\textbackslash{}int\_0\^{}1 x\^{}2\textbackslash{},dx =
\textbackslash{}frac\{1\}\{3\}</latex>

\subsection*{각주와 그림 (Footnote and figure)}\label{각주와-그림-footnote-and-figure}

각주가 붙은 문장입니다.\footnote{각주의 내용입니다.} 자동 링크 \url{https://pandoc.org} 와
\url{mailto:knix008@naver.com} 도 있습니다.

\includegraphics[width=\linewidth,keepaspectratio]{../public/app-icon.svg}

\begin{center}\rule{0.5\linewidth}{0.5pt}\end{center}

마지막 문단입니다. 줄 끝에 두 칸을 두면\\
강제 줄바꿈이 됩니다.

\end{document}
